\documentclass[11pt]{article}
% claims-sync: status=AUTHOR_REVIEW manifest=CLAIMS.md

\usepackage[T1]{fontenc}
\usepackage[utf8]{inputenc}
\usepackage{lmodern}
\usepackage{amsmath,amssymb,amsthm,mathtools}
\usepackage{booktabs}
\usepackage{geometry}
\usepackage{hyperref}
\usepackage{xcolor}

\geometry{margin=1in}

\hypersetup{
  colorlinks=true,
  linkcolor=blue!55!black,
  citecolor=blue!55!black,
  urlcolor=blue!55!black,
  unicode=true,
  pdftitle={The Regular Cyclic Point Maximizes Plücker Participation on the Positive Grassmannian},
  pdfauthor={Adam Guetz},
  pdfsubject={Extremal property of the regular cyclic point for the Plücker participation ratio, via the comass of the all-ones Plücker form},
  pdfkeywords={positive Grassmannian, Plücker coordinates, participation ratio, comass, Wirtinger inequality, regular cyclic point}
}

\newtheorem{theorem}{Theorem}
\newtheorem{lemma}[theorem]{Lemma}
\newtheorem{proposition}[theorem]{Proposition}
\newtheorem{corollary}[theorem]{Corollary}
\newtheorem{definition}[theorem]{Definition}
\newtheorem{remark}[theorem]{Remark}

\newcommand{\Gr}{\operatorname{Gr}}
\newcommand{\Part}{\mathsf{Part}}
\newcommand{\R}{\mathbb{R}}
\newcommand{\C}{\mathbb{C}}
\newcommand{\abs}[1]{\left|#1\right|}
\newcommand{\norm}[1]{\left\|#1\right\|}

\title{The Regular Cyclic Point Maximizes Plücker Participation on the Positive Grassmannian}
\author{Adam Guetz\thanks{Learned Response Labs; ORCID \href{https://orcid.org/0000-0003-0105-4628}{0000-0003-0105-4628}; \texttt{\href{mailto:learnedresponselabs@gmail.com}{learnedresponselabs@gmail.com}}. See ``A note on process'' for authorship and review status. A reading edition of this paper beside the Lean source, with narrated interactive explainers, is at \href{https://learnedresponselabs.github.io/papers/plucker_participation_cyclic_point.html}{\mbox{\nolinkurl{https://learnedresponselabs.github.io/papers/plucker_participation_cyclic_point.html}}}.}}
\date{November 24, 2026}

\begin{document}

\maketitle

\begin{abstract}
We present a proof candidate, pending external mathematical review, for the
maximum over the positive top cell \(\Gr_{>0}(k,n)\) of the
Plücker participation ratio
\[
  \Part(C)=
  \frac{\left(\sum_{|I|=k}\Delta_I(C)\right)^2}
       {\sum_{|I|=k}\Delta_I(C)^2}.
\]
The maximum is attained at the regular cyclic point and is
\[
  \begin{cases}
  \displaystyle
  \prod_{r=1}^{m}\cot^2\!\left(\frac{(2r-1)\pi}{2n}\right),
  & k=2m,\\[1.2em]
  \displaystyle
  n\prod_{r=1}^{m}\cot^2\!\left(\frac{r\pi}{n}\right),
  & k=2m+1.
  \end{cases}
\]
Moreover the regular cyclic point is the \emph{unique} global maximizer (for
each \((k,n)\) the equality locus is a single cyclic Fourier plane).
The proof converts the positive Plücker participation problem into the comass
of the all-ones Plücker \(k\)-form.  That form is
\(\Omega^m/m!\) in even degree and \(\eta\wedge\Omega^m/m!\) in odd degree,
where \(\Omega=\sum_{i<j} e_i\wedge e_j\) is the sign-kernel two-form and
\(\eta=\sum_i e_i\).  A finite Fourier diagonalization of the sign kernel,
followed by the standard Wirtinger comass bound for powers of a two-form, gives
the displayed product.  The equality planes are the cyclic Fourier planes,
whose Plücker coordinates are the regular cyclic chord products.
\end{abstract}

\paragraph{Scope.}
This note presents proof candidates for the positive top-cell bound and
uniqueness (Corollary~\ref{cor:uniqueness}) and their extension to the
nonnegative Grassmannian, with strict inequality at every boundary point
(Corollary~\ref{cor:tnn}); external mathematical review is pending. It does
not determine the maximum or maximizers on individual lower positroid cells
or boundary strata, or assert an optimization theorem for arbitrary
non-decomposable signed coordinate vectors.  The term
\emph{Plücker participation ratio} is used throughout; the same statistic is
also the inverse Simpson concentration of the normalized Plücker weights.

\noindent\begin{minipage}{\textwidth}
\tableofcontents
\end{minipage}

\section{Introduction}

The positive Grassmannian has a distinguished regular cyclic point: the point
whose Plücker coordinates are proportional to the chord products of \(n\)
regularly spaced points on a circle.  This point appears naturally in the
cyclic symmetry and critical-variety literature around the positive
Grassmannian \cite{Postnikov2006,Karp2018,GalashinCritical}.  The purpose of
this note is to record a sharp extremal property of that point for a simple
scale-invariant functional on Plücker coordinates.

Ogranovich~\cite{ogranovich2020regular} minimizes the largest-to-smallest
Plücker \emph{spread} on \(\Gr_{>0}(2,n)\), proving the regular-polygon point
extremal for \(k=2\) and stating the general-\(k\) spread extremality as a
conjecture (his Conjecture~1.2); here we instead maximize the inverse-Simpson
\emph{participation} \((\sum_I\Delta_I)^2/\sum_I\Delta_I^2\) on \(\Gr_{>0}(k,n)\)
for all \(k\), via the comass of the all-ones Plücker form, with an explicit
\(\cot^2\) value.  The two functionals differ, so our all-\(k\) proposed theorem is
complementary to --- and does not resolve --- his conjecture. The uniqueness argument appears in
Corollary~\ref{cor:uniqueness}, using the rigidity clause of the Wirtinger
lemma; it is a proof candidate pending external review.

For a positive Plücker vector \(p=(p_I)\), define
\[
  \Part(p)=\frac{\norm{p}_1^2}{\norm{p}_2^2}.
\]
Intuitively, this ratio measures how many Plücker coordinates meaningfully
contribute to \(p\): values near \(1\) indicate concentration, while large
values indicate that the Plücker mass is distributed as evenly as the Plücker
relations allow.
Equivalently, after normalizing \(w_I=p_I/\sum_J p_J\), one has
\[
  \Part(p)=\frac{1}{\sum_I w_I^2}.
\]
Thus \(\Part\) is the inverse Simpson concentration of the normalized Plücker
weights.  We use the more descriptive algebraic-geometry phrase
\emph{Plücker participation ratio}.
With \(N=\binom nk\) and \(u\) the uniform weight vector, \(u_I=1/N\),
\[
  \norm{w-u}_2^2=\sum_I w_I^2-\frac1N=\frac{1}{\Part(p)}-\frac1N .
\]
So maximizing \(\Part\) is the same as minimizing the Euclidean distance from the normalized
Plücker weights to the uniform weights.

The functional is defined using Plücker coordinates in the fixed standard
coordinate basis of \(\R^n\).  It is invariant under change of row basis and
global Plücker rescaling, hence is well-defined on \(\Gr(k,n)\), but it is not
invariant under the positive diagonal torus acting on the ambient coordinates:
torus rescaling moves the Grassmannian point and can change \(\Part\).  The
statistic therefore does not descend without further normalization to the
configuration quotient \(\Gr(k,n)/T\).  This is not an additional hypothesis
in Theorem~\ref{thm:main}; it is a boundary on how the statistic should be
reused outside this paper.

\begin{theorem}[top-cell Plücker participation maximum]
\label{thm:main}
Let \(1\leq k\leq n\), and let \(C\in\Gr_{>0}(k,n)\).  Then
\[
  \Part(C)
  =
  \frac{\left(\sum_{|I|=k}\Delta_I(C)\right)^2}
       {\sum_{|I|=k}\Delta_I(C)^2}
\]
is bounded as follows.
If \(k=2m\), then
\[
  \Part(C)
  \leq
  \prod_{r=1}^{m}\cot^2\!\left(\frac{(2r-1)\pi}{2n}\right).
\]
If \(k=2m+1\), then
\[
  \Part(C)
  \leq
  n\prod_{r=1}^{m}\cot^2\!\left(\frac{r\pi}{n}\right).
\]
Both bounds are attained by the regular cyclic point, and by no other point of
\(\Gr_{>0}(k,n)\) (Corollary~\ref{cor:uniqueness}).
\end{theorem}

Empty products are interpreted as \(1\).  Thus the formula includes, for
example, the case \(k=1\), where it reduces to
\(\Part(C)\leq n\).  At the other endpoint \(k=n\), the cotangent products
reduce to \(1\), matching the single Plücker coordinate of \(\Gr(n,n)\).

The proof actually bounds a larger problem: it computes the comass of the
all-ones Plücker form on all oriented \(k\)-planes in \(\R^n\).  Positivity
enters when identifying \(\Part(C)\) with the square of that form on the
positive Plücker representative and when verifying that the cyclic equality
plane lies in \(\Gr_{>0}(k,n)\).

The same comass computation also gives a signed decomposable variant.  If
\(E\) is an oriented \(k\)-plane with nonzero Plücker representative
\(\xi=\sum_I p_Ie_I\), define
\[
  \Part_{\pm}(E)
  =
  \frac{\left(\sum_{|I|=k}p_I\right)^2}{\sum_{|I|=k}p_I^2}.
\]
The square makes \(\Part_{\pm}\) independent of reversing the orientation of
the plane.  Unlike \(\Part\) on the positive Grassmannian, this signed
functional allows cancellations among the Plücker coordinates.

\section{Participation and the All-Ones Plücker Form}

Fix the orthonormal basis \(e_0,\ldots,e_{n-1}\) of \(\R^n\).  All subsets
\(I\) are subsets of the zero-based ground set \(\{0,\ldots,n-1\}\).  For
\(I=\{i_1<\cdots<i_k\}\), write
\[
  e_I=e_{i_1}\wedge\cdots\wedge e_{i_k}.
\]
Define the all-ones Plücker \(k\)-form
\[
  \Phi_k=\sum_{|I|=k} e_I.
\]
Here and below the standard inner product identifies \(\bigwedge^k\R^n\) with its dual, so
\(\Phi_k\) acts as a \(k\)-form: \(\Phi_k(\xi)=\langle\Phi_k,\xi\rangle=\sum_I p_I\) for
\(\xi=\sum_I p_Ie_I\).  The Plücker relations constrain the decomposable vector \(\xi\) on which
\(\Phi_k\) is evaluated.  They place no condition on \(\Phi_k\) itself, which is in general not
decomposable.
For a \(d\)-form \(\alpha\) on an inner-product space \(W\), write
\[
  \norm{\alpha}_{*,W}
  =
  \sup\{\abs{\alpha(\xi)}:\xi\in\bigwedge^d W
  \text{ is unit and decomposable}\}.
\]
When the ambient space is clear we write simply \(\norm{\alpha}_*\).

\begin{proposition}[reduction to comass]
\label{prop:comass-reduction}
Let \(C\in\Gr_{>0}(k,n)\), and let
\[
  \xi=\sum_{|I|=k}p_I e_I
\]
be a decomposable representative with all \(p_I=\Delta_I(C)>0\).  Then
\[
  \Part(C)=\Phi_k(\widehat{\xi})^2,
  \qquad
  \widehat{\xi}=\frac{\xi}{\norm{\xi}}.
\]
Consequently,
\[
  \Part(C)\leq \norm{\Phi_k}_*^2,
\]
where \(\norm{\cdot}_*\) denotes comass.
\end{proposition}

\begin{proof}
Since the exterior basis \(e_I\) is orthonormal,
\[
  \norm{\xi}^2=\sum_I p_I^2.
\]
Also
\[
  \Phi_k(\xi)=\sum_I p_I.
\]
Dividing by \(\norm{\xi}\) gives
\[
  \Phi_k(\widehat{\xi})^2
  =
  \frac{\left(\sum_I p_I\right)^2}{\sum_I p_I^2}
  =
  \Part(C).
\]
The comass bound follows because \(\widehat{\xi}\) is a unit decomposable
\(k\)-vector, so it is one of the vectors over which the comass supremum is
taken.
\end{proof}

\paragraph{Example: \(\Gr_{>0}(2,4)\).}
Order the six coordinates as \((p_{01},p_{02},p_{03},p_{12},p_{13},p_{23})\).  The coefficient
vector of \(\Phi_2\) has \(\norm{\Phi_2}^2=6\), and Cauchy--Schwarz gives \(\Part\le6\) with
equality only when all six coordinates have a common value \(t\).  That vector is not
decomposable: the Plücker relation
\[
  p_{01}p_{23}-p_{02}p_{13}+p_{03}p_{12}=0
\]
would read \(t^2=0\).  The regular cyclic point has coordinates proportional to
\((1,\sqrt2,1,1,\sqrt2,1)\), the sides and diagonals of a square.  For it the relation reads
\(1-2+1=0\), and
\[
  \Part=\frac{(4+2\sqrt2)^2}{8}=3+2\sqrt2=\cot^2\frac{\pi}{8}\approx5.83 .
\]
By Theorem~\ref{thm:main} this is the maximum.  So \(\norm{\Phi_2}^2=6\) while
\(\norm{\Phi_2}_*^2=3+2\sqrt2\): the comass is the bound that takes the Plücker relation into
account.

\section{The Sign-Kernel Two-Form}

Let
\[
  \eta=\sum_{i=0}^{n-1} e_i,
  \qquad
  \Omega=\sum_{0\leq i<j\leq n-1} e_i\wedge e_j.
\]
The key point is that the all-ones Plücker forms are powers of the same
two-form.  Here and throughout, a power of a form is its exterior power:
\(\Omega^m=\Omega\wedge\cdots\wedge\Omega\) with \(m\) factors, and \(\Omega^0=1\).

\begin{lemma}[exterior identities]
\label{lem:exterior-identities}
For \(m\geq 0\),
\[
  \Phi_{2m}=\frac{\Omega^m}{m!},
  \qquad
  \Phi_{2m+1}=\eta\wedge\frac{\Omega^m}{m!}.
\]
\end{lemma}

\begin{proof}
Let \(C\) be a skew matrix and \(\omega=\sum_{i<j}C_{ij}\,e_i\wedge e_j\). For
\(I=\{i_1<\cdots<i_{2m}\}\), the coefficient of \(e_I\) in \(\omega^m/m!\) is the
Pfaffian of the principal submatrix \(C_I\). This is the standard exterior-algebra
form of the Pfaffian; for a \(2m\times 2m\) matrix it reads
\(\omega^m/m!=\operatorname{Pf}(C)\,e_1\wedge\cdots\wedge e_{2m}\). It follows from
the definition of the Pfaffian as a signed sum over perfect matchings
\cite[\S2]{Stembridge1990}. Two-forms commute, so \(\omega^m/m!\) is the sum, over
the perfect matchings \(\pi\) of \(I\) into pairs \(i<j\), of
\(\prod_{(i,j)\in\pi}C_{ij}\) times the wedge of the \(e_i\wedge e_j\) listed by
smaller end. Sorting that wedge into \(e_I\) takes one transposition for each two
pairs of \(\pi\) that cross and zero or two for each two that do not, so its sign
is the sign of the matching.

For \(\omega=\Omega\), the matrix \(C_I\) is the \(2m\times 2m\) skew matrix
\(A_I\) with entries
\[
  (A_I)_{ab}=\operatorname{sgn}(i_b-i_a).
\]
In the ordered basis this is the skew matrix with \(1\)'s above the diagonal
and \(-1\)'s below.  Its Pfaffian is \(1\)
\cite[Proposition~2.3(c)]{Stembridge1990}.  One quick proof is induction by
expanding along the first row: every even principal deletion has the same
form, and the alternating sum \(1-1+\cdots+1\) over \(2m-1\) terms equals
\(1\).

For odd degree, let \(|I|=2m+1\). Since
\(e_{i_q}\wedge e_{I\setminus\{i_q\}}=(-1)^{q+1}e_I\), the coefficient of
\(e_I\) in \(\eta\wedge\Omega^m/m!\) is
\[
  \sum_{q=1}^{2m+1}(-1)^{q+1}
  \operatorname{Pf}(A_{I\setminus\{i_q\}}).
\]
Each even Pfaffian in this sum is \(1\), and the alternating sum has an odd
number of terms.  Hence the coefficient is again \(1\).
\end{proof}

\section{Fourier Spectrum of the Sign Kernel}\label{sec:fourier}

Let \(A\) be the skew matrix of \(\Omega\):
\[
  A_{ij}=
  \begin{cases}
    1, & i<j,\\
   -1, & i>j,\\
    0, & i=j.
  \end{cases}
\]
Throughout, the singular values of a two-form are listed once per singular pair.  Each is a
singular value of the form's skew matrix, where it occurs twice.

\begin{lemma}[spectrum on \(\R^n\)]
\label{lem:spectrum-rn}
The positive singular values of \(\Omega\) on \(\R^n\) are
\[
  \lambda_r=\cot\!\left(\frac{(2r-1)\pi}{2n}\right),
  \qquad
  r=1,\ldots,\left\lfloor\frac n2\right\rfloor,
\]
listed in decreasing order.
\end{lemma}

\begin{proof}
For a complex Fourier vector
\[
  v_z=(1,z,z^2,\ldots,z^{n-1}),
\]
if \(z^n=-1\), then a direct geometric-sum computation gives
\[
  (A v_z)_i
  =
  \sum_{j>i}z^j-\sum_{j<i}z^j
  =
  z^i\frac{1+z}{1-z}.
\]
Thus \(v_z\) is an eigenvector with eigenvalue \((1+z)/(1-z)\).  For
\(z=e^{i\theta}\),
\[
  \frac{1+z}{1-z}=i\cot(\theta/2).
\]
The roots with \(z^n=-1\) have \(\theta=(2r-1)\pi/n\), and the real skew
matrix has singular values given by the absolute values of these imaginary
eigenvalues.  Pairing conjugate roots gives the displayed positive list.
\end{proof}

The odd-degree case uses the restriction of \(\Omega\) to
\(\eta^\perp\).

\begin{lemma}[spectrum on \(\eta^\perp\)]
\label{lem:spectrum-eta-perp}
The positive singular values of \(\Omega|_{\eta^\perp}\) are
\[
  \mu_r=\cot\!\left(\frac{r\pi}{n}\right),
  \qquad
  r=1,\ldots,\left\lfloor\frac{n-1}{2}\right\rfloor,
\]
listed in decreasing order.
\end{lemma}

\begin{proof}
Let \(z^n=1\) and \(z\neq 1\).  Then \(v_z\in\eta^\perp_\C\), and
\[
  A v_z
  =
  \frac{1+z}{1-z}v_z-\frac{2}{1-z}\eta.
\]
After projecting away the \(\eta\)-component, \(v_z\) is an eigenvector of
the projected operator with eigenvalue \((1+z)/(1-z)\).  Writing
\(z=e^{2\pi i r/n}\) gives \(i\cot(r\pi/n)\), and conjugate pairing gives
the displayed singular values.
\end{proof}

\section{Comass and the Sharp Bound}

We use the standard Wirtinger comass calculation for powers of a two-form.  The
classical calibration perspective is discussed in \cite{HarveyLawson1982}.
We give the norm computation and the equality argument needed here directly.

\begin{lemma}[Wirtinger comass for powers of a two-form]
\label{lem:wirtinger}
Let \(\Theta\) be a real two-form on a Euclidean vector space with singular
values
\[
  \sigma_1\geq \sigma_2\geq \cdots \geq \sigma_s\geq 0.
\]
For \(0\le m\le s\), with the empty product understood as one,
\[
  \left\|\frac{\Theta^m}{m!}\right\|_*
  =
  \sigma_1\sigma_2\cdots\sigma_m.
\]
The supremum defining the comass is attained on the \(2m\)-plane spanned by the
first \(m\) singular pairs.  Moreover \emph{(rigidity)}, if \(m\ge1\) and
\(\sigma_m>\sigma_{m+1}\) then it is attained \emph{only} on that plane; in particular, when the top \(m\) singular
values are distinct from the remaining ones the maximizing \(2m\)-plane is
\emph{unique}.
\end{lemma}

\begin{proof}
Choose an orthonormal singular-pair basis in which
\(\Theta=\sum_{r=1}^{s}\sigma_r\,x_r\wedge y_r\). The two-forms \(x_r\wedge y_r\) commute and
have square zero, so expanding the power gives
\[
  \frac{\Theta^m}{m!}
  =\sum_{\substack{J\subseteq\{1,\dots,s\}\\ |J|=m}}
   \Bigl(\prod_{r\in J}\sigma_r\Bigr)\bigwedge_{r\in J}(x_r\wedge y_r),
\]
each \(m\)-set \(J\) arising from its \(m!\) orderings. In particular
\(\Theta^m/m!\) takes the value \(\sigma_1\cdots\sigma_m\) on the unit simple
\(2m\)-vector of the first \(m\) singular pairs, so the comass is at least this
product. Conversely, for a unit simple \(2m\)-vector
\(\xi\) spanning an oriented \(2m\)-plane \(P\), evaluation against an orthonormal
oriented basis of \(P\) gives
\(\langle\Theta^m/m!,\xi\rangle=\operatorname{Pf}\!\bigl(\Theta|_P\bigr)\), the
Pfaffian of the restriction of \(\Theta\) to \(P\), by the identity at the start of the proof of
Lemma~\ref{lem:exterior-identities} in that basis. Restriction to a subspace is a
compression of the skew operator dual to \(\Theta\). (Convention: we list a skew
form's \emph{singular pair-values} \(\beta_1\ge\beta_2\ge\cdots\), one entry per
singular pair --- the underlying matrix's singular values occur doubled.) By the
Cauchy/Poincar\'e
interlacing theorem for compressions, applied to the Hermitian complexification
used below, the \(j\)-th singular pair-value of \(\Theta|_P\)
is at most \(\sigma_j\) for every \(j\). Since \(\bigl|\operatorname{Pf}(\Theta|_P)\bigr|\)
is the product of the \(m\) singular pair-values of \(\Theta|_P\),
\[
  \bigl|\langle\Theta^m/m!,\xi\rangle\bigr|
  =\prod_{j=1}^{m}\bigl(j\text{-th singular pair-value of }\Theta|_P\bigr)
  \;\le\;\sigma_1\cdots\sigma_m,
\]
with equality on the span of the first \(m\) singular pairs, as computed above; this is
the comass value.

For the rigidity, assume \(\sigma_m>\sigma_{m+1}\) (convention:
\(\sigma_{m+1}:=0\) when \(m\) equals the number of singular pairs) and
equality. The following interlacing equality argument includes multiplicities. Let \(A\) be the real skew operator dual to
\(\Theta\) and complexify: \(H:=iA_{\C}\) is Hermitian on \(V_{\C}\), with
eigenvalue list, in decreasing order,
\(\sigma_1\ge\cdots\ge\sigma_s\ge0=\cdots=0\ge-\sigma_s\ge\cdots\ge-\sigma_1\).
Let \(P_{\C}\subset V_{\C}\) be the complexification of \(P\) and \(K\) the
compression of \(H\) to \(P_{\C}\); since compression commutes with
complexification and with multiplication by \(i\), \(K\) is Hermitian with
eigenvalues \(\pm\beta_1,\dots,\pm\beta_m\), the \(\beta_j\) being the singular
pair-values of \(\Theta|_P\). Equality gives \(\beta_j=\sigma_j\) for all
\(j\le m\).

\emph{Top family.} \(\lambda_1(K)=\sigma_1=\lambda_1(H)\), so a unit top
eigenvector \(\psi_1\in P_{\C}\) of \(K\) attains the global Rayleigh maximum
of \(H\) and is therefore an \(H\)-eigenvector with eigenvalue \(\sigma_1\).
As \(\psi_1\) is an \(H\)-eigenvector, \(\psi_1^\perp\) is \(H\)-invariant,
\(H|_{\psi_1^\perp}\) has the spectrum of \(H\) with one copy of \(\sigma_1\)
removed, and the compression of \(H|_{\psi_1^\perp}\) to
\(P_{\C}\cap\psi_1^\perp\) is \(K|_{P_{\C}\cap\psi_1^\perp}\) (for
\(v\in P_{\C}\cap\psi_1^\perp\) one has
\(\langle Hv,\psi_1\rangle=\sigma_1\langle v,\psi_1\rangle=0\)), with top
eigenvalue \(\lambda_2(K)=\sigma_2=\lambda_1(H|_{\psi_1^\perp})\). Iterating
\(m\) times produces orthonormal \(\psi_1,\dots,\psi_m\in P_{\C}\) with
\(H\psi_j=\sigma_j\psi_j\).

\emph{Bottom family.} The bottom \(m\) eigenvalues of \(K\) are
\(-\sigma_m,\dots,-\sigma_1\), equal to the bottom \(m\) of \(H\); the same
induction applied to \(-H\) produces orthonormal
\(\varphi_1,\dots,\varphi_m\in P_{\C}\) with \(H\varphi_j=-\sigma_j\varphi_j\).
Every \(\sigma_j\) with \(j\le m\) is strictly positive
(\(\sigma_j\ge\sigma_m>\sigma_{m+1}\ge0\)), so the two families have disjoint
eigenvalue sets and are mutually orthogonal: together they form \(2m\)
orthonormal vectors in the \(2m\)-dimensional \(P_{\C}\), hence span it.

\emph{Identification.} For \(j\le m\), the full \(H\)-eigenspaces of
\(\pm\sigma_j\) lie inside
\(\bigl(\bigoplus_{r\le m}\operatorname{span}(x_r,y_r)\bigr)_{\C}\): a pair
index \(r>m\) with \(\sigma_r=\sigma_j\) would force
\(\sigma_j\le\sigma_{m+1}<\sigma_m\le\sigma_j\), a contradiction. Hence
\(P_{\C}\subseteq\bigl(\bigoplus_{r\le m}\operatorname{span}(x_r,y_r)\bigr)_{\C}\);
both sides have complex dimension \(2m\), so they are equal, and taking real
points, \(P=\bigoplus_{r\le m}\operatorname{span}(x_r,y_r)\). The maximizing
plane is therefore unique.
\end{proof}

\begin{remark}[on the rigidity clause]
\label{rem:rigidity-clause}
The rigidity clause runs by Pfaffian compression and interlacing; we invoke
Harvey--Lawson only for the comass value and the calibration convention, not
for the equality-case carry-through. Its hypothesis \(\sigma_m>\sigma_{m+1}\)
is immediate for the sign-kernel cotangents (they are strictly monotone on
their ranges; \S\ref{sec:fourier}).

The rigidity proof allows ties within the selected group and a zero
next pair-value, but requires a strict gap at the cutoff. Kernel directions
cannot occur in an equality plane under that hypothesis. The wedge-line
classification requires a nonzero residual form; every application here
satisfies this requirement. Positivity and the orientation of the Fourier
witness are established by Lemma~\ref{lem:fourier-minors}.
\end{remark}

\begin{lemma}[wedging with a distinguished line]
\label{lem:wedge-line}
Use the Euclidean metric to identify vectors and covectors.
Let \(V=\R u\oplus W\) be an orthogonal decomposition with \(\norm{u}=1\).
If \(\alpha\in\bigwedge^d W^*\), then
\[
  \norm{u^*\wedge\alpha}_{*,V}=\norm{\alpha}_{*,W},
\]
with equality classification for \(\alpha\neq0\): a unit decomposable \((d+1)\)-vector attains
\(\norm{u^*\wedge\alpha}_{*,V}\) only if its plane is \(\R u\oplus Q'\), where
\(Q'\subset W\) is a \(d\)-plane whose unit \(d\)-vector attains
\(\norm{\alpha}_{*,W}\).
Consequently, for any nonzero \(\eta\in V^*\) and
\(\alpha\in\bigwedge^d(\eta^\perp)^*\),
\[
  \norm{\eta\wedge\alpha}_{*,V}
  =
  \norm{\eta}\,\norm{\alpha}_{*,\eta^\perp}.
\]
Moreover, if a two-form \(\Theta\) is decomposed as
\(\Theta=\Theta_{\eta^\perp}+\eta_0\wedge\tau\), where
\(\eta_0=\eta/\norm{\eta}\), then
\[
  \eta\wedge\Theta^m=\eta\wedge\Theta_{\eta^\perp}^m.
\]
\end{lemma}

\begin{proof}
For any unit decomposable \(d\)-vector \(\xi\in\bigwedge^d W\), the vector
\(u\wedge\xi\) is unit decomposable in \(V\), so
\[
  \norm{u^*\wedge\alpha}_{*,V}\geq\norm{\alpha}_{*,W}.
\]
Conversely, let \(\xi\) be a unit decomposable \((d+1)\)-vector with plane
\(Q\). Since \(W=u^\perp\), the subspace \(Q\cap u^\perp\) has dimension at
least \(d\); choose an orthonormal set \(q_2,\dots,q_{d+1}\) spanning a
\(d\)-dimensional subspace of \(Q\cap u^\perp\subset W\), and complete it by a
unit \(q_1\in Q\) orthogonal to it, so
\(\xi=\pm\,q_1\wedge q_2\wedge\cdots\wedge q_{d+1}\). Because
\(u^*(q_j)=0\) for \(j\ge2\), the evaluation collapses to
\[
  (u^*\wedge\alpha)(q_1,\dots,q_{d+1})
  =\langle u,q_1\rangle\,\alpha(q_2,\dots,q_{d+1}),
\]
whence
\(\bigl|(u^*\wedge\alpha)(\xi)\bigr|
\le\lvert\langle u,q_1\rangle\rvert\,\norm{\alpha}_{*,W}
\le\norm{\alpha}_{*,W}\), proving the first identity. \emph{Equality
classification:} Since a nonzero form evaluates nontrivially on some
coordinate simple wedge, \(\alpha\neq0\) implies \(\norm{\alpha}_{*,W}>0\). Thus equality forces
\(\lvert\langle u,q_1\rangle\rvert=1\), i.e.\ \(q_1=\pm u\), so
\(Q=\R u\oplus\operatorname{span}(q_2,\dots,q_{d+1})\) with
\(q_2\wedge\cdots\wedge q_{d+1}\) a unit \(d\)-vector of \(W\) attaining
\(\norm{\alpha}_{*,W}\). Scaling \(u^*\) gives the second identity and the
same classification with \(\R\eta_0\) in place of \(\R u\).  The last
identity follows because every term in \(\Theta^m\) containing the factor
\(\eta_0\wedge\tau\) vanishes after wedging with \(\eta\).
\end{proof}

\begin{proposition}[comass of the all-ones form]
\label{prop:comass-values}
If \(k=2m\), then
\[
  \norm{\Phi_{2m}}_*
  =
  \prod_{r=1}^{m}\cot\!\left(\frac{(2r-1)\pi}{2n}\right).
\]
If \(k=2m+1\), then
\[
  \norm{\Phi_{2m+1}}_*
  =
  \sqrt{n}\prod_{r=1}^{m}\cot\!\left(\frac{r\pi}{n}\right).
\]
\end{proposition}

\begin{proof}
For even \(k=2m\), combine Lemma~\ref{lem:exterior-identities},
Lemma~\ref{lem:spectrum-rn}, and Lemma~\ref{lem:wirtinger}.

For odd \(k=2m+1\), decompose
\[
  \R^n=\R\eta\oplus\eta^\perp.
\]
By Lemma~\ref{lem:wedge-line}, only the restriction of \(\Omega\) to
\(\eta^\perp\) contributes in \(\eta\wedge\Omega^m/m!\), and wedging by
\(\eta\) multiplies the comass by \(\norm{\eta}=\sqrt n\).  Therefore
Lemma~\ref{lem:spectrum-eta-perp} and Lemma~\ref{lem:wirtinger} give
\[
  \norm{\eta\wedge\Omega^m/m!}_*
  =
  \sqrt n\prod_{r=1}^{m}\cot(r\pi/n).
\]
\end{proof}

\begin{corollary}[signed participation on oriented planes]
\label{cor:signed-oriented-planes}
Let \(E\in\Gr(k,n)\) be an oriented \(k\)-plane, and let
\(\Part_{\pm}(E)\) be defined as above.  Then the same bounds as in
Theorem~\ref{thm:main} hold:
\[
  \Part_{\pm}(E)
  \leq
  \prod_{r=1}^{m}\cot^2\!\left(\frac{(2r-1)\pi}{2n}\right)
  \qquad (k=2m),
\]
and
\[
  \Part_{\pm}(E)
  \leq
  n\prod_{r=1}^{m}\cot^2\!\left(\frac{r\pi}{n}\right)
  \qquad (k=2m+1).
\]
The cyclic Fourier equality planes attain these bounds.
\end{corollary}

\begin{proof}
Let \(\xi=\sum_I p_Ie_I\) be a nonzero decomposable Plücker representative of
the plane and set \(\widehat{\xi}=\xi/\norm{\xi}\).  Then
\[
  \Part_{\pm}(E)=\Phi_k(\widehat{\xi})^2.
\]
The bound is therefore \(\norm{\Phi_k}_*^2\), whose value is computed in
Proposition~\ref{prop:comass-values}.  The equality planes are the Wirtinger
equality planes used in Section~\ref{sec:cyclic-equality-plane}; in the
positive orientation they are the regular cyclic point.
\end{proof}

\begin{proof}[Proof of Theorem~\ref{thm:main}]
Proposition~\ref{prop:comass-reduction} gives
\[
  \Part(C)\leq \norm{\Phi_k}_*^2.
\]
Squaring the values in Proposition~\ref{prop:comass-values} gives the two
displayed bounds.  Equality is treated in the next section:
Lemma~\ref{lem:fourier-minors} gives a positive equality plane, and
Corollary~\ref{cor:uniqueness} shows that it is the only one.
\end{proof}

\section{The Regular Cyclic Equality Plane}
\label{sec:cyclic-equality-plane}

The equality planes in Proposition~\ref{prop:comass-values} are explicit
Fourier subspaces.

\begin{itemize}
\item If \(k=2m\), take the real and imaginary parts of the \(m\) Fourier
  modes with frequencies \((2r-1)/(2n)\), \(r=1,\ldots,m\).
\item If \(k=2m+1\), take \(\eta/\sqrt n\) together with the real and imaginary
  parts of the \(m\) Fourier modes in \(\eta^\perp\) with frequencies
  \(r/n\), \(r=1,\ldots,m\).
\end{itemize}

Here is an explicit column identification, including orientation and scale.
Put \(\theta_j=2\pi j/n\), \(0\leq j<n\), and \(m=\lfloor k/2\rfloor\).
Define the real \(k\times n\) matrix \(K=K_{k,n}\), with rows numbered from
one, by the following ordered pairs, for \(1\leq r\leq m\):
\begin{align}
k=2m:\quad
 (K_{2r-1,j},K_{2r,j})
 &=\bigl(\cos((r-\tfrac12)\theta_j),
          \sin((r-\tfrac12)\theta_j)\bigr),\notag\\
k=2m+1:\quad K_{1,j}=1,\qquad
 (K_{2r,j},K_{2r+1,j})
 &=\bigl(\cos(r\theta_j),\sin(r\theta_j)\bigr).
\label{eq:fourier-karp-rows}
\end{align}
Thus \(K(:,j)=f_k(\theta_j)\) for the moment curve in Karp's
Equation~(1.3). In his Proposition~2.5, taking the phase parameter
\(\vartheta=-2\pi/n\) and the column index \(\ell=j+1\) gives exactly
these columns \cite[arXiv~v2 numbering]{Karp2018}. An oriented orthonormal
representative of the equality plane is
\begin{equation}
 F=D_{k,n}K,\qquad
 D_{k,n}=\begin{cases}
    \sqrt{2/n}\,I_{2m},&k=2m,\\
    \operatorname{diag}\bigl(n^{-1/2},\sqrt{2/n}\,I_{2m}\bigr),&k=2m+1.
 \end{cases}
\label{eq:fourier-karp-match}
\end{equation}
The Fourier orthogonality sums give \(FF^{\mathsf T}=I_k\), also when
\(k=n\). Every factor in \(D_{k,n}\) is positive; the identification is
a change of row basis with all ambient columns fixed.

\begin{lemma}[Fourier equality-plane minors]
\label{lem:fourier-minors}
For \(1\leq k\leq n\), the representative \(F\) in
\eqref{eq:fourier-karp-match} has, for every ordered set
\(I=\{i_1<\cdots<i_k\}\subseteq\{0,\ldots,n-1\}\),
\[
  \Delta_I(F)=n^{-k/2}
    \prod_{a<b}\abs{\zeta^{i_b}-\zeta^{i_a}}>0,
  \qquad
  \zeta=e^{2\pi i/n}.
\]
Writing \(M=\binom{k}{2}\), the corresponding unnormalized formula is
\[
 \Delta_I(K)=2^{M-m}\prod_{a<b}
      \sin\!\left(\frac{\pi(i_b-i_a)}n\right).
\]
\end{lemma}

\begin{proof}
Convert the ordered real rows of \(K\) to complex exponential rows with
increasing frequencies
\[
 -\frac{k-1}{2},-\frac{k-3}{2},\ldots,\frac{k-1}{2}.
\]
Call the resulting matrix \(V=T_kK\). Each cosine--sine pair first becomes
\((\cos-i\sin,\cos+i\sin)\), with conversion determinant \(2i\).
The permutation to increasing frequencies has \(m(m-1)\) inversions in
even rank and \(m^2\) in odd rank, including the constant row. Hence
\(\det T_k=2^m i^M\) in both cases.

On the columns \(I\), the Vandermonde identity gives
\begin{align*}
 \det V_I
 &=\prod_{a=1}^k e^{-i(k-1)\theta_{i_a}/2}
      \prod_{a<b}\bigl(e^{i\theta_{i_b}}-e^{i\theta_{i_a}}\bigr)\\
 &=(2i)^M\prod_{a<b}
      \sin\!\left(\frac{\theta_{i_b}-\theta_{i_a}}2\right).
\end{align*}
In the second line, use
\(e^{i\beta}-e^{i\alpha}
 =2i e^{i(\alpha+\beta)/2}\sin((\beta-\alpha)/2)\).
Each angle occurs \(k-1\) times in the product of differences, so the
phases cancel exactly. Dividing by \(\det T_k\) yields the stated
positive constant \(2^{M-m}\), agreeing with Karp's Lemma~2.3 because
\(M-m=\lfloor(k-1)^2/2\rfloor\).

Every sine factor is positive since \(0<i_b-i_a<n\). Replacing a sine
factor by half its chord length gives
\(\Delta_I(K)=2^{-m}\prod_{a<b}|\zeta^{i_b}-\zeta^{i_a}|\).
Finally,
\(\det D_{k,n}=(2/n)^m n^{-(k-2m)/2}\), which gives the formula for
\(F\). For \(k=1\) the products are empty and the entries are
\(n^{-1/2}\). For \(k=n\), positivity and orthonormality make the single
maximal minor exactly \(1\).
\end{proof}

The signed wrap is equally explicit. Let \(R\) be block diagonal with
the usual real rotation blocks of angles \(2\pi(r-\tfrac12)/n\) in
even rank or \(2\pi r/n\) in odd rank, and a leading \(1\) in odd rank.
The addition formulas give
\begin{equation}
 RF=[F_1\mid\cdots\mid F_{n-1}\mid(-1)^{k-1}F_0],\qquad
 \det R=1,\qquad R^n=(-1)^{k-1}I_k.
\label{eq:fourier-signed-wrap}
\end{equation}
Thus this is the positive cyclic fixed point of Karp's Theorem~1.1;
his Formula~(1.4) is the sine-product formula above with its common
projective factor suppressed \cite{Karp2018}.
Galashin's critical-variety setting supplies adjacent regular-polygon and
Ptolemy-flavored context for the same sine/chord geometry
\cite{GalashinCritical}.  The present proof uses these sources for the cyclic
equality witness, not for the aggregate cotangent product calculation.

Therefore the comass equality plane is represented by a point of
\(\Gr_{>0}(k,n)\), and the bounds in Theorem~\ref{thm:main} are sharp.

\begin{corollary}[Uniqueness of the maximizer]
\label{cor:uniqueness}
For \(k=1\), Cauchy--Schwarz gives \(\Part(C)\le n\), with equality
exactly when all positive coordinates agree, yielding the line \(\R\eta\).
For \(k\ge2\), the singular values of the sign kernel are \(\cot((2r-1)\pi/2n)\) (even degree)
and \(\cot(r\pi/n)\) (odd degree), which are \emph{strictly decreasing and
distinct} on the listed ranges, so the separation \(\sigma_m>\sigma_{m+1}\) of
Lemma~\ref{lem:wirtinger} holds.  By the rigidity clause the relevant
maximizing plane is \emph{unique}: the top \(m\) singular-pair plane in even
degree, and \(\R\eta\) plus the top \(m\) singular pairs in odd degree;
the reduction chain (Lemma~\ref{lem:exterior-identities},
Proposition~\ref{prop:comass-reduction}, Lemma~\ref{lem:wedge-line}) preserves
the equality case (for Lemma~\ref{lem:wedge-line}, by its equality
classification), and
Lemma~\ref{lem:fourier-minors} identifies that unique
plane with the regular cyclic point in \(\Gr_{>0}(k,n)\).  Hence the regular
cyclic point is the \emph{unique} maximizer of \(\Part\) on \(\Gr_{>0}(k,n)\).
\end{corollary}

\noindent The sharpness argument for Theorem~\ref{thm:main} is
independent of this corollary; the corollary upgrades sharpness to uniqueness,
its only inputs being the rigidity clause of Lemma~\ref{lem:wirtinger} (see
Remark~\ref{rem:rigidity-clause}), the strict monotonicity of
the sign-kernel cotangents, and the equality-witness identification of
Lemma~\ref{lem:fourier-minors}.

\begin{corollary}[the totally nonnegative Grassmannian]\label{cor:tnn}
Let \(C\in\Gr_{\ge0}(k,n)\), with \(\Part(C)\) defined by the same formula.  Then \(\Part(C)\)
satisfies the bounds of Theorem~\ref{thm:main}, with equality only at the regular cyclic point.
In particular, \(\Part\) is strictly below the top-cell maximum at every point of the boundary
\(\Gr_{\ge0}(k,n)\setminus\Gr_{>0}(k,n)\).
\end{corollary}

\begin{proof}
A point of \(\Gr_{\ge0}(k,n)\) has a Plücker representative with all \(p_I\ge0\), so
\(\Part(C)=\Part_{\pm}(E)\) for the underlying plane \(E\), and
Corollary~\ref{cor:signed-oriented-planes} gives the bounds.  If equality holds, then
\(\abs{\Phi_k(\widehat{\xi})}=\norm{\Phi_k}_*\), so \(E\) is a comass-maximizing plane.  The
proof of Corollary~\ref{cor:uniqueness} shows that there is only one such plane, and by
Lemma~\ref{lem:fourier-minors} it is the regular cyclic point, which lies in \(\Gr_{>0}(k,n)\).
\end{proof}

\noindent The corollary does not determine the maximum of \(\Part\) on a given boundary cell.

\section{What Is Not Claimed}

The preceding results concern the maximum of participation and its equality
locus on the totally nonnegative Grassmannian
(Theorem~\ref{thm:main}, Corollary~\ref{cor:tnn}).  Two neighboring questions
remain outside the claim boundary.

\begin{enumerate}
\item It does not prove that the same cyclic point controls every lower
  positroid cell or every boundary stratum.
\item It does not classify absolute maxima over arbitrary signed coordinate
  vectors in \(\bigwedge^k\R^n\) after dropping decomposability.  That is a
  different ambient-sphere problem: Cauchy--Schwarz gives the all-ones vector,
  which is generally not decomposable.  Corollary~\ref{cor:signed-oriented-planes}
  is only the signed decomposable, or oriented-plane, statement.
\end{enumerate}

\section*{A note on process}
Learned Response Labs is a research collaboration between the author and AI systems. The proofs, text,
software, atlases, and other artifacts of this program were primarily generated by AI coding harnesses
(OpenAI Codex and Anthropic Claude Code) operating adversarially: results produced by one system were
frozen, independently rederived, and hostilely audited by the other before internal acceptance, with
machine-checked Lean~4/Mathlib components where noted. The author designed and directed the research
program --- the choice of problems, the harness and adversarial-audit architecture, and the cross-domain
references that shaped it. \par Each paper has a \emph{reading edition}, which sets every formal statement
beside what checks it (a frozen Lean statement or a pinned exact certificate) and every step without one
beside its source. The author checks the paper's main statements, and the definitions they use, against
the edition; a version is archived and announced only after that check is complete, and the paper's page
says where the check stands. External review is pending, and reviews received will be noted with the
paper's public reading edition. \par The decision to circulate this document, and the
judgment that the process behind it makes the results worth a referee's time, are the author's; the
mathematics itself awaits external verification, which is exactly what this document requests.
Accordingly, the headline results are offered as \emph{proof candidates pending external review}, and no
claim in this document should be treated as established until that review has occurred. This
external-review preprint is prepared for public circulation; it is not yet an arXiv version.

\bibliographystyle{plain}
\bibliography{references}

\clearpage
\appendix
\begingroup\sloppy\emergencystretch=3em
\providecommand{\CyclicArchiveDOI}{DOI-PENDING}
\section*{Appendix: Formal statements and their mathematical interface}
The formal source is the \texttt{cyclic\_point/} directory of
\texttt{positive-geometry-lean}, version 0.2.0 (archive DOI: \texttt{\CyclicArchiveDOI}).
The paper remains a proof candidate pending external mathematical review.
Here ``proved'' refers to the displayed Lean statement with all its hypotheses,
not to independent validation of the translation from the written mathematics.
The project uses Lean \texttt{4.33.0-rc1} and the Mathlib revision in its lockfile.
Its 68 comparison configurations cover 193 distinct theorem names, including supporting algebra.
The checking procedure compares each solution statement against its frozen statement,
checks its axioms, and replays the exported proof through the Lean kernel and through nanoda, an independent implementation of Lean's type checker.
Only \texttt{propext}, \texttt{Classical.choice} and \texttt{Quot.sound} are admitted.
Frozen files use proof placeholders and are not themselves accepted proofs.
The following list distinguishes checked conclusions from the portions of the written argument that are not formalized.

\paragraph{Theorem~\ref{thm:main}: proved.}
For $1\le k\le n$, on matrices with positive increasing maximal minors; equality means equality of row spans. The formal headline uses a compound Rayleigh argument.
Checked statements: \texttt{CyclicPoint.\allowbreak Layer6.\allowbreak flagship}, \texttt{CyclicPoint.\allowbreak Layer6.\allowbreak karp\_\allowbreak part}, \texttt{CyclicPoint.\allowbreak Layer6.\allowbreak cyclicValue\_\allowbreak self}.

\paragraph{Proposition~\ref{prop:comass-reduction}: proved under stated hypotheses.}
Full row rank is explicit. The formal comass is defined using orthonormal row frames.
Checked statements: \texttt{CyclicPoint.\allowbreak Layer8b.\allowbreak part\_\allowbreak mul\_\allowbreak gram}, \texttt{CyclicPoint.\allowbreak Layer8b.\allowbreak part\_\allowbreak le\_\allowbreak sq\_\allowbreak comass}, \texttt{CyclicPoint.\allowbreak Layer8b.\allowbreak part\_\allowbreak eq\_\allowbreak sq\_\allowbreak comass\_\allowbreak iff}.

\paragraph{Lemma~\ref{lem:exterior-identities}: proved.}
Alternating forms are represented by their evaluation on row frames.
Checked statements: \texttt{CyclicPoint.\allowbreak Layer4.\allowbreak signMatrix\_\allowbreak pf}, \texttt{CyclicPoint.\allowbreak Layer8b.\allowbreak allOnes\_\allowbreak even\_\allowbreak eq}, \texttt{CyclicPoint.\allowbreak Layer8b.\allowbreak allOnes\_\allowbreak odd\_\allowbreak eq}.

\paragraph{Lemma~\ref{lem:spectrum-rn}: proved under stated hypotheses.}
The dimension bounds in the Fourier statements are retained.
Checked statements: \texttt{CyclicPoint.\allowbreak Layer6.\allowbreak evenFourier\_\allowbreak orthonormal}, \texttt{CyclicPoint.\allowbreak Layer6.\allowbreak evenFourier\_\allowbreak eigen}, \texttt{CyclicPoint.\allowbreak Layer6.\allowbreak evenPairs\_\allowbreak order}, \texttt{CyclicPoint.\allowbreak Layer6.\allowbreak evenFourier\_\allowbreak top}.

\paragraph{Lemma~\ref{lem:spectrum-eta-perp}: proved under stated hypotheses.}
The compression uses the orthogonal complement of the distinguished all-ones vector.
Checked statements: \texttt{CyclicPoint.\allowbreak Layer6.\allowbreak oddFourier\_\allowbreak frame}, \texttt{CyclicPoint.\allowbreak Layer6.\allowbreak oddFourier\_\allowbreak eigen}, \texttt{CyclicPoint.\allowbreak Layer6.\allowbreak oddPairs\_\allowbreak order}, \texttt{CyclicPoint.\allowbreak Layer6.\allowbreak oddFourier\_\allowbreak top}.

\paragraph{Lemma~\ref{lem:wirtinger}: proved under stated hypotheses.}
The rigidity statement retains the strict spectral gap. Its proof uses spectral trace saturation; the recursive deflation argument is not separately formalized.
Checked statements: \texttt{CyclicPoint.\allowbreak Layer8.\allowbreak skewNormalForm\_\allowbreak exists}, \texttt{CyclicPoint.\allowbreak Layer8.\allowbreak compression\_\allowbreak interlacing}, \texttt{CyclicPoint.\allowbreak Layer8.\allowbreak interlacing\_\allowbreak equality}, \texttt{CyclicPoint.\allowbreak Layer8.\allowbreak powerForm\_\allowbreak eval}, \texttt{CyclicPoint.\allowbreak Layer8.\allowbreak wirtinger\_\allowbreak value}, \texttt{CyclicPoint.\allowbreak Layer8.\allowbreak wirtinger\_\allowbreak rigid}.

\paragraph{Remark~\ref{rem:rigidity-clause}: proved under stated hypotheses.}
The strict-gap consequence and the strict ordering of the explicit pair values are separate statements.
Checked statements: \texttt{CyclicPoint.\allowbreak Layer8.\allowbreak wirtinger\_\allowbreak rigid}, \texttt{CyclicPoint.\allowbreak Layer6.\allowbreak evenPairs\_\allowbreak order}, \texttt{CyclicPoint.\allowbreak Layer6.\allowbreak oddPairs\_\allowbreak order}.

\paragraph{Lemma~\ref{lem:wedge-line}: proved under stated hypotheses.}
The equality classification requires the residual alternating form to be nonzero. Without this hypothesis a zero form attains its comass on every plane.
Checked statements: \texttt{CyclicPoint.\allowbreak Layer8b.\allowbreak wedgeOne\_\allowbreak eval}, \texttt{CyclicPoint.\allowbreak Layer8b.\allowbreak wedge\_\allowbreak line\_\allowbreak comass}, \texttt{CyclicPoint.\allowbreak Layer8b.\allowbreak wedge\_\allowbreak line\_\allowbreak rigid}, \texttt{CyclicPoint.\allowbreak Layer8b.\allowbreak wedge\_\allowbreak eta\_\allowbreak comass}, \texttt{CyclicPoint.\allowbreak Layer8b.\allowbreak wedge\_\allowbreak powerForm\_\allowbreak perp}.

\paragraph{Proposition~\ref{prop:comass-values}: proved under stated hypotheses.}
The even and odd dimension bounds and the normalization of the distinguished vector are explicit.
Checked statements: \texttt{CyclicPoint.\allowbreak Layer8b.\allowbreak comass\_\allowbreak allOnes\_\allowbreak even}, \texttt{CyclicPoint.\allowbreak Layer8b.\allowbreak comass\_\allowbreak allOnes\_\allowbreak odd}.

\paragraph{Corollary~\ref{cor:signed-oriented-planes}: proved under stated hypotheses.}
The bound and equality statement require full row rank and allow either global orientation.
Checked statements: \texttt{CyclicPoint.\allowbreak Layer6.\allowbreak flagship\_\allowbreak bound}, \texttt{CyclicPoint.\allowbreak Layer6.\allowbreak flagship\_\allowbreak equality}.

\paragraph{Lemma~\ref{lem:fourier-minors}: proved in part.}
The sine-product formula for the matrix $K$, its positive minors and existence of a determinant-one signed row shift are checked. The chord-length formula for $F$, the explicit rotation matrix and its power are not separately frozen statements.
Checked statements: \texttt{CyclicPoint.\allowbreak Layer6.\allowbreak karp\_\allowbreak even\_\allowbreak frame}, \texttt{CyclicPoint.\allowbreak Layer6.\allowbreak karp\_\allowbreak odd\_\allowbreak frame}, \texttt{CyclicPoint.\allowbreak Layer6.\allowbreak karp\_\allowbreak minor}, \texttt{CyclicPoint.\allowbreak Layer6.\allowbreak karp\_\allowbreak positive}, \texttt{CyclicPoint.\allowbreak Layer6.\allowbreak karp\_\allowbreak twistedShift}.

\paragraph{Corollary~\ref{cor:uniqueness}: proved under stated hypotheses.}
The equality plane is checked. The written comass chain is not assembled into a separate formal headline; the checked headline follows a different route.
Checked statements: \texttt{CyclicPoint.\allowbreak Layer6.\allowbreak flagship}, \texttt{CyclicPoint.\allowbreak Layer6.\allowbreak flagship\_\allowbreak equality}.

\paragraph{Corollary~\ref{cor:tnn}: proved under stated hypotheses.}
For a full-row-rank representative with nonnegative maximal minors, the bound, equality plane and strict boundary inequality are checked; this is pointwise strictness, with no uniform gap or boundary-cell extremum asserted. The supporting statement \texttt{zero\_minor\_strict}, which the paper does not state, gives strictness at a vanishing minor without a sign hypothesis.
Checked statements: \texttt{CyclicPoint.\allowbreak Layer9.\allowbreak nonneg\_\allowbreak bound\_\allowbreak equality}, \texttt{CyclicPoint.\allowbreak Layer9.\allowbreak nonneg\_\allowbreak boundary\_\allowbreak strict}, \texttt{CyclicPoint.\allowbreak Layer9.\allowbreak zero\_\allowbreak minor\_\allowbreak strict}.

\paragraph{Supporting statements and interfaces.}
The restriction identities \texttt{CyclicPoint.\allowbreak Layer8b.\allowbreak pullback\_\allowbreak eval}, \texttt{CyclicPoint.\allowbreak Layer8b.\allowbreak pullback\_\allowbreak powerForm}, \texttt{CyclicPoint.\allowbreak Layer8b.\allowbreak pullback\_\allowbreak attains}, and the determinant evaluations \texttt{CyclicPoint.\allowbreak Layer6.\allowbreak even\_\allowbreak value}, \texttt{CyclicPoint.\allowbreak Layer6.\allowbreak odd\_\allowbreak value}, supply intermediate steps.
The correspondence between row-frame evaluations and exterior algebra, the chord-length normalization, and identification with the cited geometric constructions require mathematical reading.
\par\endgroup

\end{document}
